\subsection{Performance Portability Metrics Revised (2017 -- 2021)}
\label{sec:pp_revised}

Many researchers adopted the \pp metric in their work in the following years, further underlining the necessity of such an objective performance portability metric. \todo{cite}
But the \pp metric was not only used by other researchers; they also tried highlighting potential problems with \pp and solutions, as well as improvements and extensions. 

2018: \cite{harrell_effective_2018}

\begin{itemize}
    \item not only performance portability, but productivity 
    \item critic: \pp does not penalize "heroic" optimizations efforts to increase its score (i.e., optimized code paths for specific architectures) $\rightarrow$ wouldn't be productive
\end{itemize}

\citet{harrell_effective_2018} extended the performance portability even further to an \textit{effective performance portability} that does not only include the above-discussed performance portability but also the productivity, i.e., the effort a programmer has to make to achieve said performance portability. 
This metric also penalizes frameworks that achieve good performance portability but only by drastically increasing the programmer's effort, e.g., introducing special code paths for some hardware platforms. 
However, in the remainder of this thesis, I will mainly focus on \textit{performance portability} only. \todo{Begründung?}


\newpage
\cite{sedova_high-performance_2018}

\enquote{The $PP_{MD}$ metric purports to evaluate performance portability but in practice it measures the price in performance that must be paid to make the application portable.} -- \cite{marowka_comparison_2023}

\begin{definition}[Portability \citet{sedova_high-performance_2018, mooney_bringing_1997}]
An application can be said to be portable if the cost of porting is less than the cost of re-writing:
$$DP = 1 - (\frac{C_P}{C_R}$$
where $C_P$ is the cost to port and $C_R$ the cost to rewrite the program. 
\end{definition}

\begin{itemize}
    \item $1$ $\rightarrow$ completely portable
    \item positive $\rightarrow$ porting is more profitable
    \item binary vs source code portability
\end{itemize}

\begin{definition}[MD Performance Portability~\citet{sedova_high-performance_2018}]

$$PP_{MD_c}(a, p, Q) = 
    \begin{cases}
        \frac{\abs{Q}}{\sum\limits_{i \in Q} S_i(a, p)} & \text{if }\abs{G} - \abs{Q} \neq 0\text{ and }\abs{G} - \abs{Q} \neq \abs{G} \\
        1  & \text{if }\abs{G} - \abs{Q} = \abs{G}\\
        0  & \text{if }\abs{G} - \abs{Q} = 0
    \end{cases}
$$

where $\abs{G}$ is the the cardinality of the set $G$ of algorithmic or programmatic components that significantly accelerate a best-performing application $a$, $\abs{Q}$ is the cardinality of the set $Q$ of all non-portable such accelerating components of $a$, $p$ are the parameters used in $a$, and $S$ is the speedup that this non-portable component $i$ of the application provides to the program: $S_i = \frac{T_0}{T_i}$, with $T_i$ as the time-to-solution with the use of $i$, and $T_0$ the time-to-solution without $i$.
\end{definition}

\begin{itemize}
    \item $G$: instruction-level, thread-level, process-level, or library-level
    \item $PP_{MD_c}$ not close to $\SI{100}{\percent}$: viele optimierte Teile eines Codes müssten getauscht werden mit portable Code (mit ähnlicher Performance)
    \item $PP_{MD_c}$ does not support multiple platforms
\end{itemize}

Summary of recommendations for portability (copied verbatim)
\begin{itemize}
    \item High-level programming interface with re-targetable back end that is standardized and supported by a number of both commercial and open-source initiatives.
    \item Encapsulation of machine specific sections of the program.
    \item An agreement about the worthiness of the portability effort.
    \item Modularity of the algorithm sections and the data structures, with an ability to rearrange data storage and even to change the algorithm if required.
    \item Use of only the portable subset of the standard language.
    \item A dedicated portable design effort and clear portability goals.
\end{itemize}


\newpage
\cite{yang_empirical_2018}

\begin{itemize}
    \item Roofline model as \textit{performance efficiency metric}
    \item Paper hauptsächlich wie man in "gutes" Roofline Model aufbaut
\end{itemize}

\begin{definition}[Roofline Performance Efficiency~\citet{yang_empirical_2018}]\label{def:roof_eff}
Given the observed code performance in floating point operations per second (\si{\flop\per\second}) $P_i(a, p)$, the peak \si{\flop\per\second} $F_i$ of architecture $i$, the peak bandwidth $B_i$ of architecture $i$, and the arithmetic intensity $I_i(a, p)$ of application $a$ on architecture $i$, the Roofline performance efficiency is defined as:

$$e_i(a, p) = \frac{P_i(a, p)}{\min(F_i, B_i \times I_i(a, p))}$$
\end{definition}
The denominator accounts for when the application is compute or bandwidth bound.


\newpage
\cite{yokota_beginners_2018}

\begin{itemize}
    \item tension between performance improvement and performance portability improvement
    \item \enquote{we empirically analyze the usability and usefulness of PPM, applied for OpenACC applications}
    \begin{enumerate}
        \item calculating PPM for multiple applications and platforms
        \item interpreting the results from the perspective of performance portability across devices
        \item attempting to improve the performance portability, i.e., increase PPM
    \end{enumerate}
    \item arch eff: compute-bound $\rightarrow$ operational throughput; memory-bound $\rightarrow$ bandwidth
    \item Problem arch eff: does not show whether the implementation \textbf{is actual efficient} (e.g., unnecessary memory loads, floptimization)
    \item same problem as in other studies: \pp heavily penalizes even one small performance efficiency value
    \item also points out the problem with different architecture (CPUs vs GPUs)
    \item low \pp: low efficiencies overall or big differences among platform-specific performance efficiencies
    \item how to improve low \pp?: add more platforms, use fewer platforms, improve performance efficiency
    \item \pp should be extended to include some kind of heterogeneity score
    \item trying to improve the performance efficiency $\rightarrow$ tension
    \begin{enumerate}
        \item performance \enquote{increases on all platforms}
        \item performance \enquote{increases on one platform, but decreases on another platform}
    \end{enumerate}
    \item did optimize the OpenACC code but it seems that not the baseline CUDA code $\rightarrow$ does the app eff still make sense?
    \item 
\end{itemize}

\begin{definition}[Architectural Efficiency~\citet{yokota_beginners_2018}]
    $$e\_arch_i(a, p) = 
    \begin{cases}
        \frac{OpT^{app}}{OpT^{HW}} & \text{if } OI^{app} \geq OI^{HW} \\
        \frac{BW^{app}}{BW^{HW}} & \text{otherwise}
    \end{cases}$$
    
\end{definition}


\newpage
2018: \citet{siklosi_heterogeneous_2018}


\newpage
2020: \cite{bertoni_performance_2020}

\pp insufficient (same \pp values for vastly different behavior) $\rightarrow$ also use standard deviation to accomplice \pp 

\begin{itemize}
    \item roofline schwierig zu erlangen 
    \item teilweise hier schon inconsistent
    \item durch haminic mean: gleiches \pp, trotz drastisch unterschiedlichen verhalten (einmal alle ungefähr gleich; einmal in schwacher und in sehr starker run)
    \item geben zusätzlich Standardabweichung an
\end{itemize}


\newpage
\cite{daniel_applying_2019}

$P_D$ doesn't capture performance and portability across platforms

$$P_D = \frac{\sum\limits_{i \in H} \Delta_{RMS}}{\abs{H}}$$

\begin{itemize}
    \item take problem size into account (since performance may depend on it)
    \item not a measure of productivity, but how smooth or though porting effort can be
\end{itemize}

\begin{definition}[Code Divergence~\cite{daniel_applying_2019}]\label{def:code_div}
The code divergence $\mathbf{D}$ is the average of the pairwise distances between applications in the set of codes $\mathbf{A}$:

$$D(A) = \binom{\abs{A}}{2}^{-1} \sum\limits_{\{a_i, a_j\} \subset A} d(a_i, a_j)$$

with $A \subset A_p$, where $A_p$ is the set of all applications which solve $p$. The distance $d(a_i, a_j)$ between applications $a_i$ and $a_j$ is the change in the number of LOC normalized to the smaller (in LOC) application:

$$d(a_i, a_j) = \frac{\abs{\text{LOC}(a_i) - \text{LOC}(a_j)}}{\min(\text{LOC}(a_i), \text{LOC}(a_j))}$$
    
\end{definition}

Question: can I single problem accept different input data and produce different output data? $\rightarrow$ new definition of a problem:
\begin{itemize}
    \item \enquote{a problem is a proposed question that requires an algorithm to be solved correctly}
\end{itemize}

\begin{definition}[Performance Distance~\cite{daniel_applying_2019}]\label{def:perf_dist}
Given a set of applications $\mathbf{A_p}$ that solve problem $\mathbf{p}$, let $f_i(a, p, s)$ be a measure of performance of application $\mathbf{a}$ on solving $\mathbf{p}$ on platform $\mathbf{i}$ with input size $\mathbf{s}$. 
The performance distance $\mathbf{\delta(a, \alpha}$ between applications $a$ and $\alpha$ is the relative error in performance.

$$\delta(a, \alpha) = \frac{\abs{f_i(a, p, s) - f_i(\alpha, p, s}}{f_i(\alpha, p, s)}$$

Note that $f_i(\alpha, p, s)$ is the observed result or the theoretical hardware performance. 
\end{definition}

\begin{definition}[Divergence RMS~\cite{daniel_applying_2019}]\label{def:div_rms}
The Divergence RMS $\Delta_{RMS}$ is the root mean square of performance distances between a set of input sizes of interest $S$:

$$\Delta_{RMS} = \sqrt{\frac{\sum\limits_{s \in S} \delta(a, \alpha)^2}{\abs{S}}}$$

\end{definition}

\begin{definition}[Performance Portability Divergence $\mathcal{P_D}$~\cite{daniel_applying_2019}]\label{def:pp_div}
The Performance Portability Divergence $\mathcal{P_D}$ is the arithmetic mean of RMS divergences across a set of platforms $H$:

$$\mathcal{P_D} = \frac{\sum\limits_{i \in H} \Delta_{RMS}}{\abs{H}}$$

\end{definition}

\begin{itemize}
    \item optimal $\mathcal{P_D}$ would be $\SI{0}{\percent}$
    \item penalizes large performance distances severely
    \item can be larger than $\SI{100}{\percent}$
    \item easy to reach $\SI{0}{\percent}$, than $\SI{100}{\percent}$ in \pp
\end{itemize}


\newpage
\cite{marowka_toward_2021}

\begin{itemize}
    \item Hauptproblem: harmonic mean
    \item Criterion 5 verletzt
    \item Criterion 3 (0 wenn cannot run) only a constrained because of the harmonic mean being undefined for $0$ values
    \item \pp to broad; research results are therefore unrealistic or biased
    \item geometric + harmonic mean won't work, arithmetic does 
    \item new platform, wobei das Verhältnis der Summe aller scores ($e_i$) gleich bleibt, ändert das Verhältnis von \pp
    \item Hat Probleme mit $0$ wenn nicht unterstützt
    \item harmonic mean tends to track more the lower values $\rightarrow$ in general smaller \pp than \mpp values
    \item calculate values for, e.g., CPUs and GPUs separately
    \item if we have no baseline (except a time that would be included in \pp), exclude it from \mpp
\end{itemize}

erste Definition:
\begin{definition}[Performance Portability \mpp~(\citet{marowka_comparison_2023})]\label{def:mpp}
    Given a problem $\mathbf{p}$ which is solved by an application $\mathbf{a}$ and a set of supported platforms $\mathbf{H}$ from \textbf{the same architecture class}, the performance portability metric $\overline{\pp}$ is defined as:

    \[ \mpp(a, p, H) = 
        \begin{cases}
            \frac{\sum\limits_{i \in H} e_i(a, p)}{\abs{H}} & \text{if }i\text{ is supported }\forall i \in H \\
            \textit{not applicable}                         & \text{otherwise}
        \end{cases} 
    \]

    where $\mathbf{e_i}$ is the performance efficiency of $\mathbf{a}$ solving $\mathbf{p}$ on platform $\mathbf{i \in H}$.
\end{definition}
$\rightarrow$ update of criterion 3 with \textit{not applicable}